% LSPV-4 REPORT::RANDCLOSE1::2026-10-01T14:46:00-07:00
% !TEX program = xelatex
\documentclass[11pt]{article}
\usepackage[margin=0.88in]{geometry}
\usepackage{amsmath,amssymb}
\begin{document}
\section*{LSPV-4: randomized closure attack on CRS--PPC/CCSPC}
\textbf{STATUS: OPEN.} P versus NP remains open. No unrestricted SAT lower bound is proved here.

\paragraph{Recovery and predecessor preflight.}
The malformed PCLOSE1 artifact was recovered without changing its scientific body as
\texttt{research/agents/lspv-4/2026-10-01\_1343\_LSPV-4\_PCLOSE1-R1.tex}. Readback verifies
\texttt{\string\end\{document\}} followed by
\texttt{\% END\_OF\_VERIFIED\_REPORT::LSPV-4::PCLOSE1-R1} as the final marker.
The newest LSPV-3 commit, TRIVUNI1, is substantive and contains its required header and final sentinel, so predecessor preflight passes.

\paragraph{CORPUS REVIEWED.}
I surveyed repository history/content and the governing protocol, including root \texttt{main.tex},
\texttt{AGENTS.md}, \texttt{CODEX.md}, the repo-scoped pnp-research skill, current TASK-001 through
TASK-007 surfaces, newest verified MAIN 1158-R1 and relevant historical MAIN material, LSPV-1 work
on SU05/Hadamard and coefficient normalization, LSPV-2 HADSQ1/POLAR1 and earlier stacked-code work,
LSPV-3 DEPDEF1/DEPSWAP1/TRIVUNI1 and reconstruction work, LSPV-4 ACQ1/REREAD1/PCLOSE1-R1, governance,
and the HEC/VSC/G-PPC/CRS-PPC/CCSPC/RCCN/Search-G causal branches. Repository search was run for
CCSPC, CRS-PPC, ACQ1, HADSQ1, PCLOSE1, current tasks, and current MAIN snapshots.

\paragraph{Two obstacles and primary-literature preflight.}
The two load-bearing obstacles attacked here are: (i) whether fresh randomized challenges can rescue
CCSPC after PCLOSE1's deterministic polynomial-depth closure; and (ii) whether any resulting argument
escapes relativization/algebrization. Primary literature checked includes Shaltiel--Umans on resilient
AM reconstruction, van Melkebeek--Mocelin Sdroievski on targeted HSGs/PCPPs, the
Sipser--G\'acs--Lautemann containment of BPP in the second level of PH, and Aaronson--Wigderson on
algebrization. Current pointer-chasing lower bounds remain communication-model results and do not
supply a centralized static-SAT runtime transfer.

\paragraph{PROVED/EXACT: randomized polynomial-depth closure under the counterfactual.}
Consider a uniform CCSPC/CRS--PPC static computation with polynomially many stages. At each stage allow
polynomial deterministic work, polynomially many adaptive queries to an NP language, and polynomially
many fresh random bits. Suppose the whole computation has bounded two-sided error. Then its language
lies in
\[
\mathrm{BPP}^{\mathrm{NP}}.
\]
Under the counterfactual \(P=NP\), every NP-oracle call can be replaced by a deterministic polynomial-time
algorithm, so the same computation lies in BPP. The Sipser--G\'acs--Lautemann theorem gives
\[
\mathrm{BPP}\subseteq \Sigma_2^P\cap\Pi_2^P\subseteq PH.
\]
Under \(P=NP\), the polynomial hierarchy collapses to P. Therefore
\[
\boxed{P=NP\Longrightarrow \mathrm{BPP}^{\mathrm{NP}}=P}
\]
for the declared polynomial-time oracle computation, and in particular the randomized polynomial-depth
CCSPC static language is in P.

The inclusion needed for the hostile conclusion is the upper bound
\(\mathrm{BPP}^{\mathrm{NP}}\subseteq P\) under \(P=NP\); the reverse inclusion is trivial. This theorem
does not provide an efficient deterministic simulation unconditionally.

\paragraph{Consequence for fresh-challenge arguments.}
Fresh post-result randomness can be essential to an online information theorem, but if the final
staticized decision problem remains a uniform bounded-error polynomial-time computation with NP
subroutines, freshness does not produce a superpolynomial resource under the \(P=NP\) counterfactual.
Thus the generic implication
\[
\boxed{\text{fresh random challenges + polynomial causal depth}
\Longrightarrow \text{superpolynomial static SAT work}}
\]
is \textbf{FALSIFIED/RETIRED}. Randomness can still participate in a lower-bound proof only if a new
model-preserving theorem forces a resource outside this P-closed simulation, rather than merely
preserving unpredictability during the online execution.

\paragraph{Guards: what this does not prove.}
This does not say every individual random seed has a deterministic polynomial-time correctness
guarantee. It uses the language-level bounded-error BPP containment. It also does not cover a
superpolynomial number of realized stages encoded succinctly at unit cost; that would require a
different input/model analysis and risks a succinct-input/NEXP confusion. Nor does it convert
communication, streaming, cell-probe, or VLSI lower bounds into centralized RAM time lower bounds.

\paragraph{Interaction with PCLOSE1 and hostile controls.}
PCLOSE1 proved the deterministic closure baseline for polynomially many polynomial-time SAT-dependent
stages. RANDCLOSE1 closes the obvious randomized escape hatch at the class level. Static reread,
pointer jumping, memoization, group compression, stagewise nondeterminism, and INTERP/RECOMP remain
additional bypasses. DEAD/COPY/DELAY can inflate transcript depth without changing the extensional
decision problem. Horn, 2-SAT, bounded-treewidth, and affine/Tseitin controls continue to show that
long causal descriptions can coexist with efficient global solution methods.

\paragraph{SU/vMS, acquisition, and reconstruction branches.}
This result does not alter the source-sanctioned SU05 Hadamard algebra. HADSQ1/POLAR1 remain exact
fixed-realization algebraic statements. The exact SU09/vMS invocation trace, global stacked Schur
rank/common kernel, coefficient-compatible sparse-\(Q\) normalization, and the exact legal-state
separation evaluator remain OPEN. ACQ1 remains CONDITIONAL on the model-matched evaluator/compiler
cost. TRIVUNI1 sharpens reconstruction: under the assumed SAT decider the trivial downstream
uniformizer \(U_A=1-A\circ D_j\) already exists at the charged composition cost, so a useful
uniformizer must quantitatively beat that baseline rather than merely exist.

\paragraph{PQ and cross-domain classification.}
\textbf{EXACT OBJECT/THEOREM TRANSFER:} fixed-code Schur products, quadratic-relation kernels,
puncturing/typed shortening, POLAR1's alternating-form equations, collision Cauchy--Schwarz, and
Sipser--G\'acs--Lautemann at the complexity-class level.
\textbf{ANALOGY/CANDIDATE:} syzygies, Hilbert functions, low-rank relation matrices, LWE-style external
extraction, special-soundness/collision extraction, communication pointer chasing, and cell-probe
lower bounds. Each still needs a SAME\_REALIZATION succinct algorithm or model-preserving transfer.
\textbf{FAILED TRANSFER:} cryptographic hardness/random-oracle assumptions, VLSI locality, distributed
cuts, online freshness by itself, and communication cost as an unrestricted centralized SAT runtime
lower bound.

\paragraph{Relativization and algebrization audit.}
The new closure theorem is not a nonrelativizing hinge. The stagewise simulation is explicitly
compatible with replacing semantic subproblems by oracle answers, and the active Schur/POLAR
machinery is algebraic. Aaronson--Wigderson therefore remains a live barrier. A successful final route
must identify a proof line whose validity is destroyed by semantic/algebraic oracle replacement while
remaining valid for every ordinary polynomial-time SAT algorithm.

\paragraph{DUPLICATION / CONTRADICTION AUDIT.}
CCSPC1 and TASK-005 already warned that a causal NDTM may wait for current challenges and that latent
branch count is not runtime. PCLOSE1 already proved deterministic polynomial-depth P-closure.
RANDCLOSE1's new exact contribution is the bounded-error randomized extension through
\(\mathrm{BPP}^{NP}\) and Sipser--G\'acs--Lautemann. No contradiction with HADSQ1, POLAR1, ACQ1,
TRIVUNI1, or the online HEC/CRS-PPC information theorems was found because those results charge
different declared resources/models.

\paragraph{STALE OR OVERSTATED MAIN CLAIMS; adopted/rejected/retired.}
Any root-MAIN reading that treats fresh CRS-PPC/CCSPC randomness as a possible standalone escape from
PCLOSE1 is now stale and should be downgraded. Adopted: source-sanctioned Hadamard branch, charged
ACQ1, TRIVUNI1's trivial-uniformizer baseline, and online causal theorems in their declared models.
Rejected as insufficient: fresh randomness, canonicalization, latent path count, and polynomial causal
depth as standalone static hardness. Retired: direct online-freshness-to-static-runtime transfer.

\paragraph{Components compatible with hypothetical \(P=NP\).}
Static reread, polynomially many adaptive SAT decisions, canonical SAT self-reduction, polynomial
randomized causal cascades, stagewise nondeterminism, memoization, group compression, and ACQ1 NP
search are all compatible with the counterfactual. This is not evidence for \(P=NP\); it is a hostile
control any contradiction must survive.

\paragraph{Strict status and next falsification target.}
\textbf{PROVED/EXACT:} under \(P=NP\), the declared uniform bounded-error polynomial-time CCSPC with
polynomially many adaptive NP calls is in P; PCLOSE1-R1 recovery is verified.
\textbf{KNOWN THEOREM:} BPP is contained in \(\Sigma_2^P\cap\Pi_2^P\).
\textbf{CONDITIONAL:} ACQ1's charged acquisition ledger.
\textbf{CANDIDATE:} a resource outside the P-closed randomized static simulation and a succinct
structure-specific algebraic separator.
\textbf{FALSIFIED/RETIRED:} fresh randomness plus polynomial causal depth as a standalone
superpolynomial static-SAT resource.
\textbf{OPEN:} exact SU09/vMS invocation, stacked Schur rank/common kernel, sparse-\(Q\) normalization,
reconstruction exponent reduction beyond the trivial \(A\circ D_j\) baseline, arbitrary-algorithm
transfer, and a nonrelativizing/nonalgebrizing hinge. Hence P versus NP remains open.

The next falsification target is the strongest proposed CRS-PPC/CCSPC static compiler: classify its
end-to-end language after all random challenges and adaptive SAT stages are included. If it remains in
\(\mathrm{BPP}^{NP}\) with polynomial realized depth, then under \(P=NP\) it is in P and cannot itself
supply the contradiction. Any surviving compiler must identify an independently proved resource not
erased by this closure, static reread, composition, memoization, or semantic/algebraic oracle access.

\begin{thebibliography}{5}
\bibitem{Sipser83} M. Sipser, \emph{A complexity theoretic approach to randomness}, STOC 1983.
\bibitem{Lautemann83} C. Lautemann, \emph{BPP and the polynomial hierarchy}, Information Processing Letters 17, 1983.
\bibitem{SU09} R. Shaltiel and C. Umans, \emph{Low-End Uniform Hardness versus Randomness Tradeoffs for AM}, SIAM J. Comput. 39(3):1006--1037, 2009.
\bibitem{vMS25} D. van Melkebeek and N. Mocelin Sdroievski, \emph{Instance-Wise Hardness and Refutation versus Derandomization for Arthur--Merlin Protocols}, Computational Complexity 34:15, 2025.
\bibitem{AW09} S. Aaronson and A. Wigderson, \emph{Algebrization: A New Barrier in Complexity Theory}, ACM Trans. Comput. Theory 1(1), 2009.
\end{thebibliography}
\end{document}
% END_OF_VERIFIED_REPORT::LSPV-4::RANDCLOSE1